\begin{definition}[Directed insertion at an arbitrary torus position]%
    \label{def:peps_torus_translated_flux_assignment}
    \lean{TNLean.PEPS.torusTranslatedFluxAssignment,
        TNLean.PEPS.torusTranslatedFluxAssignment_eq_treeCycleAssignment,
        TNLean.PEPS.torusTranslatedFluxAssignment_up_transport}
    \leanok
    \uses{thm:peps_translated_two_plaquette_geometry,
        thm:peps_regular_two_cycle_flux_move}
    On a finite torus of horizontal period at least four and vertical period
    at least three, fix any position $v=(x,y)$. Let $u_g^0$ give upward
    transport $g$ on the vertical bond at $v$, and let $u_g^1$ give upward
    transport $g$ on both vertical bonds at $v$ and $v+(1,0)$.
    Every other bond carries the identity. In the ordered-edge convention,
    both selected bonds receive the same element
    \begin{align}
        k_y(g)=
        \begin{cases}
            g,&\operatorname{val}(y)<\operatorname{val}(y+1),\\
            g^{-1},&\text{otherwise}.
        \end{cases}
    \end{align}
    The literal assignments are precisely the tree-cycle assignments of the
    translated six-site block. This specifies the insertion in
    \cite[Theorem~6.16]{Schuch2010PEPS} at either periodic seam as well.
\end{definition}

\begin{theorem}[The actual flux endpoint moves at every position]%
    \label{thm:peps_torus_translated_flux_holonomy}
    \lean{TNLean.PEPS.regularWalkHolonomy_torusTranslatedFluxAssignment}
    \leanok
    \uses{def:peps_torus_translated_flux_assignment,
        def:peps_torus_plaquette_flux_geometry}
    For the counterclockwise native plaquette walks based at $v$ and $v+(1,0)$,
    the actual holonomies satisfy
    \begin{align}
        (H_L(u_g^0),H_R(u_g^0))&=(g^{-1},1),&
        (H_L(u_g^1),H_R(u_g^1))&=(1,g^{-1}).
    \end{align}
    These formulas hold at every position, including both periodic seams.
    Thus the insertion removes the flux at the left plaquette and places the
    same flux at the right plaquette, as in
    \cite[Theorem~6.16]{Schuch2010PEPS}.
\end{theorem}

\begin{proof}\leanok
    The directed upward transports are $g$ on the selected bonds, regardless
    of their native endpoint sorting. The left plaquette traverses the left
    insertion downward and the middle insertion upward, so its extended
    holonomy is $g^{-1}g=1$. The right plaquette traverses the middle insertion
    downward, giving $g^{-1}$. Its other vertical bond is distinct from both
    selected bonds under the period bounds. Every horizontal transport is
    the identity. Retaining only the left insertion gives the initial values.
\end{proof}

\begin{theorem}[An actual six-spin flux operation at every position]%
    \label{thm:peps_torus_translated_physical_flux_move}
    \lean{TNLean.PEPS.IsGIsometric.exists_unitary_torusTranslatedPhysicalFluxMove}
    \leanok
    \uses{thm:peps_translated_two_plaquette_geometry,
        def:peps_torus_translated_flux_assignment,
        thm:peps_regular_two_cycle_physical_flux_move}
    Let the original four-leg tensor be $G$-isometric for the regular virtual
    action. For every position $v$, there exists a unitary $W_v$ on the six
    original spins of the cyclic rectangle $R_v$ such that
    \begin{align}
        W_v\Psi_{R_v}(u_g^0,\theta)=\Psi_{R_v}(u_g^1,\theta)
    \end{align}
    for every $g\in G$ and every crossing configuration $\theta$.
    Both vectors are the actual open-region contractions of the original
    tensor. The unitary is chosen before $g$ and $\theta$; no spanning tree,
    Gram identity, state equality, or state translation invariance is supplied.
\end{theorem}

\begin{proof}\leanok
    Derive the actual spanning path, root, and distinct non-tree chords from
    the translated numbering. Native four-leg $G$-isometry gives the required
    incident-edge isometry at every site. Apply the fixed physical two-cycle
    operation to these derived data, with its common insertion element
    $k_y(g)$. The literal assignment identities give the desired actual columns.
\end{proof}

\begin{theorem}[A global torus flux movement at every position]%
    \label{thm:peps_torus_translated_global_flux_move}
    \lean{TNLean.PEPS.IsGIsometric.exists_unitary_torusTranslatedGlobalFluxMove}
    \leanok
    \uses{thm:peps_translated_two_plaquette_geometry,
        def:peps_torus_translated_flux_assignment,
        thm:peps_torus_translated_flux_holonomy,
        thm:peps_regular_two_cycle_global_flux_move}
    Under the same period and local $G$-isometry hypotheses, for every
    position $v$ there is one unitary $W_v$ on the six original spins such
    that its complementary identity extension $\widehat W_v$ is a global
    unitary and
    \begin{align}
        \widehat W_v\Psi(u_g^0;u)=\Psi(u_g^1;u)
    \end{align}
    for every $g\in G$ and every common assignment $u$ on exterior and
    crossing bonds. The internal bonds have the specified literal insertions;
    the other internal bonds carry the identity. Both states are actual
    globally contracted torus states. The unitary precedes $g$ and $u$.
    This is the arbitrary-position regular finite-torus realization of the
    local flux movement in \cite[Theorem~6.16]{Schuch2010PEPS}, including either
    periodic seam. The remaining unrestricted lattice scope is recorded in
    \cite{gap:rmp_peps_quantum_double_g_isometry}.
\end{theorem}

\begin{proof}\leanok
    Apply the global physical two-cycle operation with the derived translated
    tree, root, and pair of chords. Both native ordered insertion operators
    equal $k_y(g)$, so the literal assignments give its single- and double-bond
    inputs. The global contraction identity preserves every common exterior
    and crossing operator. The complementary identity extension preserves
    unitarity. The proved plaquette holonomies give the stated flux movement.
\end{proof}
